\section{Análisis}

Mismo firmware que \texttt{resultados/2026-10-05\_pll} ($F_{cy} = 50$\,MHz,
con reporte de la causa de reinicio); la tarjeta no se reprogramó. El único
cambio es la conexión USB: el FT2232H y la Curiosity Nano, que antes iban en
puertos distintos de la PC, ahora comparten un hub conectado a un solo
puerto.

\begin{table}[H]
\centering
\begin{tabular}{@{}lrr@{}}
\toprule
 & Puertos distintos & Un solo hub \\
\midrule
CRC incorrecto en la PC        & 0.68\,\% & 0 \\
Número de secuencia            & 0.73\,\% & 0 \\
Longitud incorrecta (dsPIC)    & 0.72\,\% & 0 \\
Reinicios reportados           & 0        & 0 \\
\bottomrule
\end{tabular}
\end{table}

Las fallas en las fronteras del byte 1 y del byte 7, que se atribuían a
\texttt{armar\_dma()}, desaparecieron sin modificar el firmware. Como dañaban
los dos sentidos de la misma transferencia a la vez, el resultado es
consistente con ruido en SCK o en \texttt{CS} por una diferencia en la
referencia de tierra entre las dos tarjetas cuando cada una se alimenta de un
puerto distinto. Sin el analizador lógico no se puede confirmar el
mecanismo; una corrida con las tarjetas de nuevo en puertos distintos
indicaría si la conexión es la causa.

La única bandera de vigilancia corresponde a la pausa entre la prueba
anterior y ésta.
